\documentclass[aspectratio=169, 11pt]{beamer}
\usepackage{beamer_preamble}

\begin{document}

% ==== Title ===============================================================
\titleslide{Lesson 6-4 \textperiodcentered\ Quick}{Logarithmic Functions + DOK-3 Sales Revenue Inverse (Item 28)}

% ==== Do Now ==============================================================
\begin{frame}{Do Now --- Log translations + inverse setup}{Algebra 2 \textperiodcentered\ Lesson 6-4}
\phasebadges{1--2}{7}{bridge to log inverses}
\vspace{0.2cm}
\begin{projcard}{Three warm-up items (solo, 7 min)}{slidegold}
\begin{enumerate}
  \setlength{\itemsep}{3pt}
  \item \textbf{q31} --- Translation multi-select for \(h(x) = \ln(x+2) - 1\). Which statements are true?
  \item \textbf{q5} --- Graph \(y = \log_4 x\). State domain, range, VA, \(x\)-intercept.
  \item \textbf{q32 (SAT/ACT)} --- Find the inverse of \(f(x) = 5^{x+1}\). (Choices A--D.)
\end{enumerate}
\end{projcard}
\end{frame}

% ==== Launch ==============================================================
\begin{frame}{Launch --- Example 3B: inverse of a log function}{Algebra 2 \textperiodcentered\ Lesson 6-4}
\phasebadges{2}{8}{teacher models}
\vspace{0.2cm}
\begin{projcard}{Rules Card --- Log Key Features + Inverse Rules (keep on your page)}{slidegreen}
\small
\textbf{Key features of \(y = \log_b x\):} Domain \(x > 0\); VA: \(x = 0\); \(x\)-int: \((1,0)\); passes through \((b, 1)\).\\[2pt]
\textbf{Translated form \(y = \log_b(x-h)+k\):} VA shifts to \(x = h\); domain \(x > h\).\\[2pt]
\textbf{Inverse rules:}\\
\hspace*{1em}\(\bullet\;\) \(y = a\cdot b^x \;\Rightarrow\; f^{-1}(x) = \log_b(x/a)\)\\
\hspace*{1em}\(\bullet\;\) \(y = \log_b(x-h)+k \;\Rightarrow\; f^{-1}(x) = b^{(x-k)}+h\)\\[2pt]
\textbf{Composition check:} \(f(f^{-1}(x)) = x\)
\end{projcard}
\end{frame}

\begin{frame}{Launch --- Example 3B walkthrough}{Algebra 2 \textperiodcentered\ Lesson 6-4}
\phasebadges{2}{8}{teacher models}
\vspace{0.2cm}
\begin{projcard}{Ex 3B: Find the inverse of \(g(x) = \log_7(x+5)\)}{slidegreen}
\small
\begin{enumerate}
  \setlength{\itemsep}{2pt}
  \item Replace \(g(x)\) with \(y\): \quad \(y = \log_7(x+5)\)
  \item Swap \(x\) and \(y\): \quad \(x = \log_7(y+5)\)
  \item Rewrite in exponential form: \quad \(7^x = y+5\)
  \item Solve for \(y\): \quad \(y = 7^x - 5\) \quad\(\Rightarrow\;\) \(g^{-1}(x) = 7^x - 5\)
  \item \textbf{Composition check:} \(g(g^{-1}(x)) = \log_7(7^x - 5 + 5) = \log_7 7^x = x\;\checkmark\)
\end{enumerate}
\vspace{3pt}
\textcolor{slidegray}{\small 90-sec contrast --- Ex 3A: inverse of \(y = 3\cdot 5^x\): same swap-and-solve, direction reversed \(\rightarrow\) \(y = \log_5(x/3)\).}
\end{projcard}
\end{frame}

% ==== Explore =============================================================
\begin{frame}{Explore --- Think-Pair-Share (25 min)}{Algebra 2 \textperiodcentered\ Lesson 6-4}
\phasebadges{2--3}{25}{student work}
\small\textcolor{slidegray}{5 items in order. Plan \(\sim\)10 min for Practice \#28 (DOK-3). 45-min Wed F: skip q9.}

\vspace{0.1cm}
\begin{projcard}{Practice \#28 \textperiodcentered\ \(\bigstar\) DOK-3 SALES REVENUE INVERSE}{slideaccent}
\small Sales revenue model. Find the inverse. Judge which equation form is easier for a given \(R\).
\end{projcard}
\vspace{-0.12cm}
\begin{projcard}{Practice \#21 \textperiodcentered\ Inverse of \(f(x) = 5^{x-3}\)}{slideblue}
\small Find the inverse. [Form B Q7 rehearsal]
\end{projcard}
\vspace{-0.12cm}
\begin{projcard}{Practice \#24 \textperiodcentered\ Inverse of \(f(x) = \log_2(8x)\)}{slideblue}
\small Swap and rewrite. Watch for domain. [Form B Q12 rehearsal]
\end{projcard}
\vspace{-0.12cm}
\begin{projcard}{Practice \#13 \textperiodcentered\ Equation from translated log graph}{slideblue}
\small Identify VA shift \(\rightarrow\) write equation. [Form B Q11 rehearsal]
\end{projcard}
\vspace{-0.12cm}
\begin{projcard}{Practice \#9 \textperiodcentered\ Are graphs inverses? [Form B Q9/Q10]}{slideblue}
\small Visual test: reflection symmetry over \(y = x\).
\end{projcard}
\end{frame}

% ==== Share / Summary =====================================================
\begin{frame}{Share / Summary}{Algebra 2 \textperiodcentered\ Lesson 6-4}
\phasebadges{2}{5}{DOK-3 reveal + 6-5 bridge}
\vspace{0.2cm}
\begin{projcard}{Sentence frame \textperiodcentered\ Practice \#28}{slideblue}
\small ``For Practice \#28, I found the inverse by \underline{\hspace{2.2cm}}.\\
For revenue \(R = \underline{\hspace{0.8cm}}\) dollars, the original form is easier because \underline{\hspace{2cm}}.\\
The inverse form is easier when \underline{\hspace{2.8cm}}.''
\end{projcard}

\vspace{0.15cm}
\begin{projcard}{Bridge to Lesson 6-5}{slideaccent}
\small Log properties (product, quotient, power rules) let you expand or condense log expressions --- making inverses and equations faster.
\end{projcard}
\end{frame}

% ==== Exit Ticket =========================================================
\begin{frame}{Exit Ticket --- Inverse + Asymptote}{Algebra 2 \textperiodcentered\ Lesson 6-4}
\phasebadges{1--2}{5}{not graded}
\vspace{0.2cm}
\begin{projcard}{Complete on your exit ticket}{slidegold}
\begin{enumerate}
  \setlength{\itemsep}{0.25cm}
  \item Find the inverse of \(f(x) = 3\cdot\log_5(x-2)+1\).
  \item State the vertical asymptote of \(f\).
  \item State the horizontal asymptote of \(f^{-1}\).
  \item One biggest thing I learned today: \underline{\hspace{4cm}}.
\end{enumerate}
\end{projcard}
\end{frame}

\end{document}
